\documentclass[10pt,a4paper]{article}
\usepackage[margin=2.2cm]{geometry}
\usepackage{fontspec}
\usepackage{kotex}
\setmainfont{Apple SD Gothic Neo}
\setsansfont{Apple SD Gothic Neo}
\usepackage{amsmath,amssymb}
\usepackage{graphicx}
\usepackage{booktabs}
\usepackage{caption}
\usepackage[dvipsnames]{xcolor}
\usepackage[colorlinks=true,linkcolor=MidnightBlue,citecolor=MidnightBlue,
            urlcolor=MidnightBlue]{hyperref}
\usepackage{titlesec}
\titleformat{\section}{\large\bfseries}{\thesection}{0.6em}{}
\titleformat{\subsection}{\normalsize\bfseries}{\thesubsection}{0.5em}{}
\captionsetup{font=small,labelfont=bf}
\setlength{\parskip}{0.35em}
\setlength{\parindent}{0pt}

% 주장 상자 — 모든 주장은 반증 조건을 함께 적는다
\usepackage{mdframed}
\newmdenv[linewidth=0.4pt,linecolor=gray!60,backgroundcolor=gray!5,
          innertopmargin=6pt,innerbottommargin=6pt,skipabove=8pt,skipbelow=8pt]{claimbox}
\newcommand{\claim}[2]{\begin{claimbox}\textbf{주장.} #1\par\smallskip
  \textcolor{gray!70}{\textbf{반증 조건.} #2}\end{claimbox}}

\title{\bfseries 상관 1.000 인 두 열이 같은 열이 아니다\\[2pt]
\large 겹말의 공유 폭}
\author{월드모델 실험실 · 노트 653--654}
\date{2026-08-05}
\begin{document}\maketitle

\section{한 줄}

순위 상관이 \textbf{정확히 1.000} 인 두 열이 있었다. 하나를 빼니 그 도메인이
$-0.0160$ 을 잃었다. 같은 조건의 다른 쌍은 빼도 $+0.0009$ 로 꿈쩍도 안 했다.
\textbf{상관은 두 열이 일치한다고만 말한다. 하나를 빼도 되는지는 말하지 않는다.}

\section{먼저, 가드가 유령을 보고 있었다}

겹말 가드(노트 240)는 축 쌍의 순위 상관을 보고 $|\rho| \ge 0.85$ 를 겹말로
적는다. 이 가드가 챔피언에서 후보 하나를 냈다 --- \verb|trend_level| 과
\verb|trend_volatility| 가 만화 $+0.998$ $\cdot$ 펀딩 $+0.992$ 다.

원인을 ``검색량이 곱셈적이라 표준편차가 평균에 비례한다''로 적고 변동계수로
고치려 했는데, 코드를 열어 보니 \verb|ingest/naver_trend.py| 가 \textbf{이미}
\verb|std/mean| 이었다. 진단이 틀렸다.

진짜 원인은 \textbf{결측}이다. 두 축이 같은 원천이라 \verb|zero_is_data=True|
가 \emph{같은 행}을 같은 상수로 채우고, 백분위 변환은 동률에 같은 값을 주므로
그 덩어리가 두 축에서 \textbf{완전히 일치}한다.

\begin{center}\small
\begin{tabular}{lrrrr}\toprule
도메인 & 채움값 몫 & $\rho$ 전체 & $\rho$ 덩어리 제외 \\ \midrule
게임 & 0.013 & $+0.010$ & $+0.007$ \\
만화 & 0.937 & $+0.998$ & $\mathbf{-0.654}$ \\
펀딩 & 0.880 & $+0.992$ & $\mathbf{-0.505}$ \\
세계애니 & 0.714 & $+0.950$ & $-0.358$ \\ \bottomrule
\end{tabular}\end{center}

\textbf{덩어리가 없으면 상관도 없다}(게임). 덩어리를 빼면 부호가 \emph{음수}로
뒤집힌다 --- 많이 찾는 것은 안정적으로 찾는다. 두 축은 실제로 다른 것을 나른다.

가드를 고쳤다(최빈값 덩어리를 벗기고 잰다). 회귀 시험: 고침을 끄면 3쌍,
켜면 \textbf{2쌍}.

\section{그리고 내 규칙이 시야를 지웠다}

후보를 고르는 규칙을 ``$|\rho| \ge 0.85$ 를 \textbf{세 도메인 이상}에서''로
못박았다. 사후 고르기를 막으려던 것인데, 그것이 \textbf{한 도메인에만 있는 축을
원리적으로 지웠다}(\verb|fund_*| $\cdot$ \verb|mob_*| $\cdot$ \verb|tag_*|).
가드는 도메인별로 훑으므로 그 규칙에 안 걸리고, $|\rho| = \mathbf{1.000}$ 인
진짜 겹말 둘을 찾았다.

\claim{사후 고르기를 막는 규칙은 시야도 같이 지운다.}{규칙에 ``무엇을 막는가''
와 함께 ``무엇을 못 보게 하는가''를 적지 않으면 그 사각을 가드가 대신 볼
때까지 모른다.}

\section{순수 중복을 뺐더니 갈렸다}

\begin{center}\small
\begin{tabular}{lrrr}\toprule
뺀 축 & 짝 & 그 도메인 & 판 \\ \midrule
\verb|fund_cat| & \verb|grp| (7도메인 공유) & $\mathbf{-0.0160}$ & $-0.0029$ (0/12) \\
\verb|mob_nlang| & \verb|target_breadth| (도메인 전용) & $+0.0009$ & $-0.0007$ (5/12) \\ \bottomrule
\end{tabular}\end{center}

둘 다 $|\rho| = 1.000$ 이고 둘 다 한 도메인만 쓰는 축이다. 뺀 것이 원래
있었는지도 확인했다(펀딩 2,887행 $\cdot$ 모바일 2,000행, 뺀 뒤 둘 다 완전히
사라짐). \textbf{그런데 한쪽은 크게 잃고 한쪽은 안 잃는다.}

\section{기제 --- 공유 폭}

노트 241 이 겹말의 이득을 ``자료가 아니라 도메인별 기울기''로 갈랐다. 그
자유가 값을 가지려면 \textbf{짝이 여러 도메인에 걸쳐 있어야} 한다.

\verb|fund_cat| 은 \verb|grp| 의 사본이고 \verb|grp| 은 \textbf{일곱 도메인}이
함께 쓴다. 사본이 있으면 펀딩은 \emph{제 열로 제 기울기}를 갖는다. 빼면 여섯
이웃과 눈금 하나를 나눠 써야 한다 --- \textbf{자유를 잃는다}.

\verb|mob_nlang| 이 겹치는 \verb|target_breadth| 는 \textbf{기반 축}이라 모바일이
이미 제 열을 갖고 있다. 잃을 자유가 없다.

\claim{$|\rho| = 1.000$ 이라도 열을 빼면 손해일 수 있고, 손해의 크기는 짝의
\textbf{공유 폭}이 정한다.}{공유 폭이 같은 두 쌍에서 제거 손해가 크게 다르면
이 설명이 부족한 것이다.}

\section{처분}

\textbf{둘 다 그대로 둔다.} 판 차가 문턱 $0.0045$ 안이고, 뺄 이유가 없다.
겹말 가드는 노트 240 의 처분대로 \textbf{거부권이 아니라 이름표}로 남는다 ---
그 처분이 다시 옳았다. 다만 이름표에 상관과 함께 \textbf{짝이 몇 도메인에
걸쳐 있나}를 같이 적어야 한다.

한 가지가 더 남는다. 펀딩 $-0.0160$ 은 0/12 이고 그 도메인 잡음을 넘는데
\textbf{판 평균이 그것을 지웠다}($-0.0029$). 노트 611 이 ``판이 다섯 자 중
제일 약하다''고 적은 자리가 여기서 또 나온다.

\begin{figure}[h]\centering
\includegraphics[width=\textwidth]{figs/f1.pdf}
\caption{(a) 채움값 덩어리의 몫이 유령 상관의 크기를 정한다 --- 덩어리를 빼면
부호가 뒤집힌다. (b) $|\rho|=1.000$ 인 두 쌍인데 제거 손해가 갈린다(회색 띠는
판 문턱 $\pm0.0045$). (c) 갈리는 이유는 짝의 공유 폭이다.}
\end{figure}

\end{document}
